\section{Password-protected files and credentials}

Not every secret needs a hardware quorum. \code{access password} protects a
file with a passphrase alone, and \code{vault} stores ordinary credentials
(a username plus a secret) the same way.

\begin{lstlisting}
keyquorum access password --state 0 --source ./secret.txt \
  --encrypted-path ./secret.txt.kqenc
keyquorum access password --state 0 --source ./secret.txt \
  --encrypted-path ./secret.txt.kqenc --expires "2026-12-31 23:59"
keyquorum access password --state 1 --id 1 --output ./secret.txt

keyquorum vault add "Email" --username alice
keyquorum vault get 1
\end{lstlisting}

Either can also take \flag{--pin} to require a 4-digit PIN, attempt-limited
in the style of a FIDO2 PIN, alongside the password. A successful one-time
PIN check is cached for one hour; end that window early with, for example:

\begin{lstlisting}
keyquorum pin relock --resource credential --id 1
\end{lstlisting}

\begin{noteBox}
Password-protected files share the same UTC \flag{--expires} cutoff and
destructive purge-on-first-touch behavior as quorum-protected files
(Section~\ref{sec:quorum-files}).
\end{noteBox}
